\documentclass[11pt]{article}
\usepackage[a4paper,margin=30mm]{geometry}
\usepackage[english]{babel}
\usepackage{fontspec}
\usepackage{microtype}
\usepackage{amsmath,amssymb,mathtools}
\usepackage{array,booktabs}
\setmainfont{DejaVu Serif}
\setsansfont{DejaVu Sans}
\setmonofont{DejaVu Sans Mono}
\setlength{\parindent}{1.2em}
\setlength{\parskip}{0.35em}
\emergencystretch=3em
\allowdisplaybreaks
\newcommand{\OO}{\mathcal O}
\newcommand{\res}{\operatorname{Res}}
\begin{document}

\section*{3.2. Singular Deformations}

According to 2.7, the versal moduli space of deformations of a singularity of a supercurve is a subspace of the versal moduli space of the underlying superscheme. To first order, one has the following result.

\paragraph{Lemma 3.3.}
The versal moduli space of a singularity of a supercurve is, to first order, the following subspace of the versal moduli space, computed in 3.1, of the superscheme singularity:
\[
\begin{array}{ll}
\text{\rm I:} & \tau=0,\\[2mm]
\text{\rm II:} & \tau_1=\tau_2=0,\qquad t_1+t_2=0 .
\end{array}
\]

The point is to know when one can deform $\delta$, while respecting the following conditions:
\begin{enumerate}
\item[(a)] $\delta$ takes values in $\omega_X$;
\item[(b)] on the smooth locus, $\delta$ satisfies 1.4(ii).
\end{enumerate}

In the deformations under consideration, parametrized by a vector subspace of the one that parametrizes the deformation of 3.1, the complement $U$ of the singular point is the union of two branches. On these branches we take the respective coordinates $(z_1,\theta_1)$ and $(z_2,\theta_2)$. We describe a section $\omega$ of $\omega_X$ by its restrictions $\omega_1$ and $\omega_2$ to the two branches. The condition on $\omega_1$ and $\omega_2$ is that, for every section $f$ of $\OO_X$, with restrictions $f_1$ and $f_2$ to the two branches, one has in $\OO_S$
\[
  \res_0(f_1\omega_1)+\res_0(f_2\omega_2)=0 .
\]
Since $\omega_X$ is flat over $S$, to obtain a set of generators of $\omega_X$ it is enough to lift a set of generators of $\omega_{X_0}$, where $X_0$ is the special fibre.

In case I, respectively II, let $b$ denote the following section of $\omega_X$ on $U$:
\[
\begin{array}{ll}
\text{\rm I:}&
\begin{cases}
(1)\ \left[\dfrac{dz_1/z_1}{d\theta}\right],\\[2mm]
(2)\ \left[-\,\dfrac{dz_2/z_2}{d\theta}\right],
\end{cases}
\\[6mm]
\text{\rm II:}&
\begin{cases}
(1)\ \left[\dfrac{dz_1}{d\theta_1}\right],\\[2mm]
(2)\ \left[\dfrac{dz_2}{d\theta_2}\right].
\end{cases}
\end{array}
\]
Let $D_0$ be the odd derivation on $U$:
\[
\begin{array}{ll}
\text{\rm I:}&
\begin{cases}
(1)\ \partial_\theta+\theta\,\partial_{z_1},\\
(2)\ \partial_\theta-\theta\,\partial_{z_2},
\end{cases}
\\[5mm]
\text{\rm II:}&
\begin{cases}
(1)\ \partial_{\theta_1}+\theta_1\,\partial_{z_1},\\
(2)\ \partial_{\theta_2}+\theta_2\,\partial_{z_2}.
\end{cases}
\end{array}
\]
Set $\delta_0=D_0\cdot b$. This $\delta_0$ does not necessarily send $\OO_X$ to $\omega_X$; for a deformation as a supercurve to be possible, one must be able to correct $\delta_0$ so that it sends $\OO_X$ to $\omega_X$.

We distinguish cases I and II, and the even and odd deformations.

\paragraph{Case I, odd.}
The deformation is
\[
  z_1z_2=\tau\theta .
\]
Here $b$ lies in $\omega_X$; one has
\[
  b=\frac{\left[dz_1\,dz_2/d\theta\right]}{d(z_1z_2-\tau\theta)},
\]
so it is a basis of $\omega_X$. One would have to correct the odd derivation $D_0$, that is, pass to $D_0+\tau D$, where $D$ is an even derivation with values in $\OO$. Describing a section of $\OO_X$ by its restrictions to the branches, one has
\[
z_1=
\begin{cases}
(1)\ z_1,\\
(2)\ \tau\theta/z_2
\end{cases}
\quad \xrightarrow{\ D_0\ }\quad
\begin{cases}
\theta z_1,\\
\tau/z_2 ,
\end{cases}
\]
which is not in $\OO_X$: $\theta z_1$ is in $\OO_X$, while $\tau/z_2$ is not, since $0$ is not in $(1/z_2)$. The derivation $D$ would have to send $z_1$ to a nonzero quantity on the second branch; this is impossible.

\paragraph{Case II. Absence of odd deformations.}
The $X_0$ is bihomogeneous; one can assign the weights
\[
  \theta_1:(1,0),\quad z_1:(2,0),\qquad
  \theta_2:(0,1),\quad z_2:(0,2).
\]
The versal deformation of the singularity is also bihomogeneous:
\[
  \tau_1:(1,2),\qquad \tau_2:(2,1),\qquad
  t_1,t_2:(1,1).
\]
Therefore, if there is a deformation with
\[
  \tau_1=a\tau,\qquad \tau_2=b\tau,\qquad (a,b)\ne(0,0),
\]
there is also one with $(a,b)$ replaced by $(\lambda a,\mu b)$. If $a$ or $b$ is zero, then by symmetry there is a deformation with $\tau_2=0$; otherwise, there are two distinct deformations and $\tau_1,\tau_2$ are in fact free.

It is therefore enough to study the deformation over $k[\tau_1]$:
\[
  z_1z_2=\tau_1\theta_1,\qquad
  z_1\theta_2=z_2\theta_1=\theta_1\theta_2=0.
\]

We compute $\omega_X$. Write an element of $\omega_X$ according to its values on each branch:
\[
  \begin{cases}
  u_1\left[\dfrac{dz_1}{d\theta_1}\right],\\[2mm]
  u_2\left[\dfrac{dz_2}{d\theta_2}\right].
  \end{cases}
\]
One has to lift the generators. The generator $\left[dz_1/d\theta_1\right]$ is not, as it stands, in $\omega_X$, because
\[
 z_2\left[\frac{dz_1}{d\theta_1}\right]
 =
 \begin{cases}
 \dfrac{\tau_1\theta_1}{z_1}\left[\dfrac{dz_1}{d\theta_1}\right],\\[2mm]
 z_2\cdot0,
 \end{cases}
\]
has residue $\tau_1$. It lifts to
\[
  \begin{cases}
  \left[\dfrac{dz_1}{d\theta_1}\right],\\[2mm]
  -\,\dfrac{\tau_1\theta_2}{z_2}\left[\dfrac{dz_2}{d\theta_2}\right].
  \end{cases}
\]
Similarly, $\left[dz_2/d\theta_2\right]$ lifts to
\[
  \begin{cases}
  0,\\[1mm]
  \left[\dfrac{dz_2}{d\theta_2}\right].
  \end{cases}
\]
Finally,
\[
  \frac{\theta_1}{z_1}\left[\frac{dz_1}{d\theta_1}\right]
  -
  \frac{\theta_2}{z_2}\left[\frac{dz_2}{d\theta_2}\right]
\]
lifts by the same expression on both branches.

Now seek to lift
\[
  \delta=D_0\cdot b
\]
as in 3.1, II. The rational lift given by the same formula on each branch is not valued in $\omega_X$: one has
\[
  \theta_1=
  \begin{cases}
  \theta_1,\\
  0
  \end{cases}
  \longmapsto
  \left[\frac{dz_1}{d\theta_1}\right].
\]
The derivation $D_0$ would have to be corrected by an even rational derivation sending
\[
  \theta_1\longmapsto
  \frac{\theta_2}{z_2}\left[\frac{dz_2}{d\theta_2}\right]
\]
plus a regular section on the second branch, which is impossible.

\paragraph{Case II, even deformations.}
One has
\[
\begin{array}{rlrl}
z_1z_2&=0, & \theta_1\theta_2&=0,\\
z_1\theta_2&=t_1\theta_1, & z_2\theta_1&=t_2\theta_2 .
\end{array}
\]
On the two branches:
\[
\begin{array}{lll}
\text{branch 1:}& z_2=0,\quad \theta_2=t_1\theta_1/z_1,
& \text{coordinates }(z_1,\theta_1),\\[1mm]
\text{branch 2:}& z_1=0,\quad \theta_1=t_2\theta_2/z_2,
& \text{coordinates }(z_2,\theta_2).
\end{array}
\]
We again look for generators of $\omega_X$, in the form
\[
\begin{cases}
u_1(z_1,\theta_1)\left[\dfrac{dz_1}{d\theta_1}\right],\\[2mm]
u_2(z_2,\theta_2)\left[\dfrac{dz_2}{d\theta_2}\right],
\end{cases}
\qquad
u_i\in k[z_i,z_i^{-1};\theta_i].
\]
The generator $\left[dz_1/d\theta_1\right]$ lifts as
\[
\begin{cases}
\left[\dfrac{dz_1}{d\theta_1}\right],\\[2mm]
-\dfrac{t_1}{z_2}\left[\dfrac{dz_2}{d\theta_2}\right],
\end{cases}
\]
and $\left[dz_2/d\theta_2\right]$ lifts as
\[
\begin{cases}
-\dfrac{t_2}{z_1}\left[\dfrac{dz_1}{d\theta_1}\right],\\[2mm]
\left[\dfrac{dz_2}{d\theta_2}\right].
\end{cases}
\]
Moreover,
\[
\begin{cases}
\dfrac{\theta_1}{z_1}\left[\dfrac{dz_1}{d\theta_1}\right],\\[2mm]
-\dfrac{\theta_2}{z_2}\left[\dfrac{dz_2}{d\theta_2}\right]
\end{cases}
\]
lifts in the same way.

Let $D_0$ be as above. One has
\[
D_0(z_1)=
\begin{cases}
\theta_1\left[\dfrac{dz_1}{d\theta_1}\right],\\
0,
\end{cases}
\qquad
D_0(z_2)=
\begin{cases}
0,\\
\theta_2\left[\dfrac{dz_2}{d\theta_2}\right],
\end{cases}
\]
which are in $\omega_X$. Next
\[
D_0(\theta_1)=
\begin{cases}
\left[\dfrac{dz_1}{d\theta_1}\right],\\[2mm]
\dfrac{t_2}{z_2}\left[\dfrac{dz_2}{d\theta_2}\right],
\end{cases}
\quad\text{is in }\omega_X\text{ if }t_1=-t_2,
\]
and
\[
D_0(\theta_2)=
\begin{cases}
\dfrac{t_1}{z_1}\left[\dfrac{dz_1}{d\theta_1}\right],\\[2mm]
\left[\dfrac{dz_2}{d\theta_2}\right],
\end{cases}
\quad\text{is in }\omega_X\text{ if }t_1=-t_2.
\]
If $t_1+t_2\ne0$, one must correct $D_0$ to $D_0+\varepsilon D$, with $D$ even; once again, $D$ would have to send $\theta_1$ to a nonzero quantity on the second branch, which is impossible.

Thus one can deform only when
\[
  t_1+t_2=0.
\]
That the non-excluded deformations are possible follows from 2.5.

\end{document}